\documentclass[a4paper,12pt]{article}
\usepackage[utf8]{inputenc}
\usepackage[ngerman]{babel}
\usepackage{amsmath,amssymb}
\usepackage{geometry}\geometry{margin=2.5cm}
\usepackage{enumitem}
\usepackage{booktabs}
\usepackage{tikz}
\usepackage{pgfplots}\pgfplotsset{compat=1.18}
\usepackage{tcolorbox}
\tcbuselibrary{breakable}
\usepackage{xcolor}
\newtcolorbox{begriffe}[1][]{colback=yellow!10, colframe=yellow!60!black, fonttitle=\bfseries, title={Was bedeuten die Begriffe?}, breakable, #1}
\newtcolorbox{wastun}[1][]{colback=orange!5, colframe=orange!60!black, fonttitle=\bfseries, title={Was ist zu tun?}, breakable, #1}
\newcommand{\piv}[1]{{\color{red}\boxed{#1}}}

\title{\"Ubungsblatt 10 -- L\"osungen \\ \large Linear Algebra}
\author{}
\date{}

\begin{document}
\maketitle

\begin{begriffe}
\begin{itemize}[leftmargin=*, nosep]
    \item \textbf{Koordinatenvektor $[v]_B$}: ``Wie viel von jedem Basisvektor brauche ich, um $v$ zusammenzubauen?'' Wenn $v = \alpha b_1 + \beta b_2 + \dots$, dann ist $[v]_B = (\alpha, \beta, \dots)$.
    \item \textbf{Darstellungsmatrix $[T]_B$}: Die Spalten sind die Bilder der Basisvektoren, ausgedr\"uckt \emph{in der Basis $B$}. Spalte $j$ $= [T(b_j)]_B$.
    \item \textbf{Standardbasis $E$}: $e_1 = (1,0,\dots)$, $e_2 = (0,1,\dots)$, \dots In dieser Basis ist $[v]_E = v$ (man liest die Koordinaten direkt ab).
    \item \textbf{Basiswechselmatrix $P$}: Spalten $=$ die Basisvektoren von $B$ in Standard-Koordinaten. Dann gilt $[v]_E = P\,[v]_B$ und umgekehrt $[v]_B = P^{-1} [v]_E$.
    \item \textbf{Basiswechsel f\"ur Abbildungen}: $[T]_E = P\,[T]_B\,P^{-1}$ \quad bzw. \quad $[T]_B = P^{-1} [T]_E\, P$.
\end{itemize}
\end{begriffe}

%% ============================================================
\subsection*{Aufgabe 39}
%% ============================================================

\textbf{Angabe.} $T \colon \mathbb{R}^2 \to \mathbb{R}^2$, $T(x, y) = (x + 2y,\; -3x)$. Basis $B = \{(0,1),\; (-1,2)\}$, Vektor $v = (1, 2)$.

Wir benennen die Basisvektoren: $b_1 = (0, 1)$, $b_2 = (-1, 2)$.

\begin{wastun}
\begin{itemize}[leftmargin=*, nosep]
    \item \textbf{(a)} $[T]_B$ auf zwei Wegen: (1) direkt -- Bilder $T(b_1), T(b_2)$ in $B$-Koordinaten als Spalten; (2) erst $[T]_E$, dann Basiswechsel mit $P$.
    \item \textbf{(b)} $[T(v)]_B$ auf zwei Wegen: (1) erst $[v]_B$, dann mit $[T]_B [v]_B$; (2) erst $T(v)$ in Standard, dann Basiswechsel.
\end{itemize}
\end{wastun}

%% --- 39 (a) ---
\subsubsection*{(a) $[T]_B$ berechnen}

\textbf{Weg 1: Direkt aus der Definition.}

\textbf{Schritt 1: Bilder der Basisvektoren ausrechnen.}
\begin{align*}
T(b_1) = T(0, 1) &= (0 + 2\cdot 1,\; -3 \cdot 0) = (2, 0) \\
T(b_2) = T(-1, 2) &= (-1 + 2\cdot 2,\; -3 \cdot (-1)) = (3, 3)
\end{align*}

\textbf{Schritt 2: Jedes Bild in $B$-Koordinaten umrechnen.} Das hei\ss t: schreibe $(p,q) = \alpha b_1 + \beta b_2 = \alpha(0,1) + \beta(-1,2) = (-\beta,\; \alpha + 2\beta)$. Vergleich liefert allgemein $\beta = -p$ und $\alpha = q + 2p$ \dots wir rechnen es einfach pro Vektor:

\textbf{(a) $T(b_1) = (2, 0)$:} $-\beta = 2 \Rightarrow \beta = -2$; \; $\alpha + 2\beta = 0 \Rightarrow \alpha = 4$.\\
Probe: $4(0,1) + (-2)(-1,2) = (2, 0)$ \checkmark. \quad Also $[T(b_1)]_B = (4, -2)$.

\textbf{(b) $T(b_2) = (3, 3)$:} $-\beta = 3 \Rightarrow \beta = -3$; \; $\alpha + 2\beta = 3 \Rightarrow \alpha = 9$.\\
Probe: $9(0,1) + (-3)(-1,2) = (3, 3)$ \checkmark. \quad Also $[T(b_2)]_B = (9, -3)$.

\textbf{Schritt 3: Spalten zusammensetzen.}
\[
[T]_B = \begin{pmatrix} 4 & 9 \\ -2 & -3 \end{pmatrix}.
\]

\bigskip

\textbf{Weg 2: Erst $[T]_E$, dann Basiswechsel.}

\textbf{Schritt 1: $[T]_E$ aufstellen} (Spalten $= T(e_1), T(e_2)$ in Standard):
\[
T(e_1) = T(1,0) = (1, -3), \quad T(e_2) = T(0,1) = (2, 0)
\;\Rightarrow\;
[T]_E = \begin{pmatrix} 1 & 2 \\ -3 & 0 \end{pmatrix}.
\]

\textbf{Schritt 2: Basiswechselmatrix $P$ und ihre Inverse.} Spalten von $P$ $=$ die $B$-Vektoren:
\[
P = \begin{pmatrix} 0 & -1 \\ 1 & 2 \end{pmatrix}, \quad
\det P = 0\cdot 2 - (-1)\cdot 1 = 1, \quad
P^{-1} = \begin{pmatrix} 2 & 1 \\ -1 & 0 \end{pmatrix}.
\]

\textbf{Schritt 3: $[T]_B = P^{-1} [T]_E\, P$ ausrechnen.} Erst $[T]_E P$:
\[
[T]_E P = \begin{pmatrix} 1 & 2 \\ -3 & 0 \end{pmatrix}\begin{pmatrix} 0 & -1 \\ 1 & 2 \end{pmatrix}
= \begin{pmatrix} 2 & 3 \\ 0 & 3 \end{pmatrix}.
\]
Dann $P^{-1}$ davor:
\[
[T]_B = \begin{pmatrix} 2 & 1 \\ -1 & 0 \end{pmatrix}\begin{pmatrix} 2 & 3 \\ 0 & 3 \end{pmatrix}
= \begin{pmatrix} 4 & 9 \\ -2 & -3 \end{pmatrix}.
\]

\[
\fbox{\parbox{0.93\textwidth}{\centering
Beide Wege: $[T]_B = \begin{pmatrix} 4 & 9 \\ -2 & -3 \end{pmatrix}$.
}}
\]

%% --- 39 (b) ---
\subsubsection*{(b) $[T(v)]_B$ berechnen, $v = (1, 2)$}

\textbf{Weg 1: Erst $[v]_B$, dann $[T]_B [v]_B$.}

\textbf{Schritt 1: $v$ in $B$-Koordinaten.} $(1, 2) = \alpha(0,1) + \beta(-1,2) = (-\beta, \alpha + 2\beta)$:
\[
-\beta = 1 \Rightarrow \beta = -1; \quad \alpha + 2\beta = 2 \Rightarrow \alpha = 4.
\]
Probe: $4(0,1) + (-1)(-1,2) = (1, 2)$ \checkmark. \quad Also $[v]_B = (4, -1)$.

\textbf{Schritt 2: Mit dem Satz $[T(v)]_B = [T]_B [v]_B$.}
\[
[T(v)]_B = \begin{pmatrix} 4 & 9 \\ -2 & -3 \end{pmatrix}\begin{pmatrix} 4 \\ -1 \end{pmatrix}
= \begin{pmatrix} 16 - 9 \\ -8 + 3 \end{pmatrix}
= \begin{pmatrix} 7 \\ -5 \end{pmatrix}.
\]

\bigskip

\textbf{Weg 2: Erst $T(v)$ in Standard, dann Basiswechsel.}

\textbf{Schritt 1: $T(v)$ direkt.} $T(1,2) = (1 + 2\cdot 2,\; -3\cdot 1) = (5, -3)$, also $[T(v)]_E = (5, -3)$.

\textbf{Schritt 2: In $B$-Koordinaten mit $P^{-1}$.}
\[
[T(v)]_B = P^{-1} [T(v)]_E = \begin{pmatrix} 2 & 1 \\ -1 & 0 \end{pmatrix}\begin{pmatrix} 5 \\ -3 \end{pmatrix}
= \begin{pmatrix} 10 - 3 \\ -5 + 0 \end{pmatrix}
= \begin{pmatrix} 7 \\ -5 \end{pmatrix}.
\]

\[
\fbox{\parbox{0.93\textwidth}{\centering
Beide Wege: $[T(v)]_B = (7, -5)$.
}}
\]

\bigskip
\hrule
\bigskip

%% ============================================================
\subsection*{Aufgabe 40}
%% ============================================================

\textbf{Angabe.} $T \colon \mathbb{R}^2 \to \mathbb{R}^2$ mit $[T]_B = \begin{pmatrix} 1 & 2 \\ 0 & -1 \end{pmatrix}$, Basis $B = \{(1,2),\; (1,3)\}$. Gesucht: explizite Formel f\"ur $T(x, y)$.

\begin{wastun}
\begin{itemize}[leftmargin=*, nosep]
    \item Wir kennen $[T]_B$, wollen aber $T$ in normalen Koordinaten. Plan: $[T]_E = P\,[T]_B\,P^{-1}$ ausrechnen ($P$ $=$ $B$-Vektoren als Spalten). Aus $[T]_E$ liest man die Formel direkt ab.
\end{itemize}
\end{wastun}

\textbf{Schritt 1: Basiswechselmatrix $P$ und Inverse.}
\[
P = \begin{pmatrix} 1 & 1 \\ 2 & 3 \end{pmatrix}, \quad
\det P = 3 - 2 = 1, \quad
P^{-1} = \begin{pmatrix} 3 & -1 \\ -2 & 1 \end{pmatrix}.
\]

\textbf{Schritt 2: $[T]_E = P\,[T]_B\,P^{-1}$ ausrechnen.} Erst $[T]_B P^{-1}$:
\[
\begin{pmatrix} 1 & 2 \\ 0 & -1 \end{pmatrix}\begin{pmatrix} 3 & -1 \\ -2 & 1 \end{pmatrix}
= \begin{pmatrix} 3 - 4 & -1 + 2 \\ 0 + 2 & 0 - 1 \end{pmatrix}
= \begin{pmatrix} -1 & 1 \\ 2 & -1 \end{pmatrix}.
\]
Dann $P$ davor:
\[
[T]_E = \begin{pmatrix} 1 & 1 \\ 2 & 3 \end{pmatrix}\begin{pmatrix} -1 & 1 \\ 2 & -1 \end{pmatrix}
= \begin{pmatrix} -1 + 2 & 1 - 1 \\ -2 + 6 & 2 - 3 \end{pmatrix}
= \begin{pmatrix} 1 & 0 \\ 4 & -1 \end{pmatrix}.
\]

\textbf{Schritt 3: Formel ablesen.} Die Matrix $[T]_E$ angewandt auf $(x, y)$:
\[
\begin{pmatrix} 1 & 0 \\ 4 & -1 \end{pmatrix}\begin{pmatrix} x \\ y \end{pmatrix}
= \begin{pmatrix} x \\ 4x - y \end{pmatrix}.
\]

\[
\fbox{\parbox{0.93\textwidth}{\centering
$T(x, y) = (x,\; 4x - y)$.
}}
\]

\textbf{Probe:} $T(1,2) = (1, 4-2) = (1, 2) = b_1$ und $T(1,3) = (1, 4-3) = (1,1)$.
In $B$-Koordinaten: $b_1 = (1,0)_B$ (1.\ Spalte von $[T]_B$ \checkmark) und $(1,1) = 2 b_1 - b_2$, also $(2,-1)_B$ (2.\ Spalte \checkmark).

\bigskip
\hrule
\bigskip

%% ============================================================
\subsection*{Aufgabe 41}
%% ============================================================

\textbf{Angabe.} $T \colon \mathbb{R}^3 \to \mathbb{R}^2$, $T(x, y, z) = (x - 2z,\; y - 3z)$.
Basen $B = \{(1,0,0),(1,1,0),(1,1,1)\}$ (im Definitionsbereich) und $C = \{(1,0),(1,1)\}$ (im Zielraum).

Wir benennen: $b_1 = (1,0,0)$, $b_2 = (1,1,0)$, $b_3 = (1,1,1)$, \; $c_1 = (1,0)$, $c_2 = (1,1)$.

\begin{begriffe}
\begin{itemize}[leftmargin=*, nosep]
    \item $[T]_B^C$: Abbildung von Basis $B$ (Quelle) nach Basis $C$ (Ziel). Spalte $j$ $= [T(b_j)]_C$.
    \item Hier sind \textbf{zwei} verschiedene Basen im Spiel -- eine f\"ur die Quelle, eine f\"ur das Ziel.
\end{itemize}
\end{begriffe}

%% --- 41 (a) ---
\subsubsection*{(a) $[T]_B^C$ berechnen}

\textbf{Weg 1: Direkt aus der Definition.}

\textbf{Schritt 1: Bilder der $B$-Vektoren.}
\begin{align*}
T(b_1) = T(1,0,0) &= (1 - 0,\; 0 - 0) = (1, 0) \\
T(b_2) = T(1,1,0) &= (1 - 0,\; 1 - 0) = (1, 1) \\
T(b_3) = T(1,1,1) &= (1 - 2,\; 1 - 3) = (-1, -2)
\end{align*}

\textbf{Schritt 2: Jedes Bild in $C$-Koordinaten.} Schreibe $(p, q) = \alpha c_1 + \beta c_2 = \alpha(1,0) + \beta(1,1) = (\alpha + \beta,\; \beta)$. Daraus: $\beta = q$, $\alpha = p - q$.

\begin{itemize}[leftmargin=*, nosep]
    \item $T(b_1) = (1,0)$: $\beta = 0$, $\alpha = 1$ \; $\Rightarrow [T(b_1)]_C = (1, 0)$.
    \item $T(b_2) = (1,1)$: $\beta = 1$, $\alpha = 0$ \; $\Rightarrow [T(b_2)]_C = (0, 1)$.
    \item $T(b_3) = (-1,-2)$: $\beta = -2$, $\alpha = -1-(-2) = 1$ \; $\Rightarrow [T(b_3)]_C = (1, -2)$.
\end{itemize}

\textbf{Schritt 3: Spalten zusammensetzen.}
\[
[T]_B^C = \begin{pmatrix} 1 & 0 & 1 \\ 0 & 1 & -2 \end{pmatrix}.
\]

\bigskip

\textbf{Weg 2: Erst Standardbasen, dann Basiswechsel.}

\textbf{Schritt 1: Matrix $A$ von $T$ in Standardbasen} (Spalten $= T(e_1), T(e_2), T(e_3)$):
\[
T(e_1) = (1, 0),\; T(e_2) = (0, 1),\; T(e_3) = (-2, -3)
\;\Rightarrow\;
A = \begin{pmatrix} 1 & 0 & -2 \\ 0 & 1 & -3 \end{pmatrix}.
\]

\textbf{Schritt 2: Basiswechselmatrizen.}
\[
P_B = \begin{pmatrix} 1 & 1 & 1 \\ 0 & 1 & 1 \\ 0 & 0 & 1 \end{pmatrix} \;(\text{$B$-Vektoren als Spalten}), \quad
P_C = \begin{pmatrix} 1 & 1 \\ 0 & 1 \end{pmatrix}, \quad
P_C^{-1} = \begin{pmatrix} 1 & -1 \\ 0 & 1 \end{pmatrix}.
\]

\textbf{Schritt 3: Formel $[T]_B^C = P_C^{-1}\, A\, P_B$.} Erst $A P_B$:
\[
A P_B = \begin{pmatrix} 1 & 0 & -2 \\ 0 & 1 & -3 \end{pmatrix}\begin{pmatrix} 1 & 1 & 1 \\ 0 & 1 & 1 \\ 0 & 0 & 1 \end{pmatrix}
= \begin{pmatrix} 1 & 1 & -1 \\ 0 & 1 & -2 \end{pmatrix}.
\]
Dann $P_C^{-1}$ davor:
\[
[T]_B^C = \begin{pmatrix} 1 & -1 \\ 0 & 1 \end{pmatrix}\begin{pmatrix} 1 & 1 & -1 \\ 0 & 1 & -2 \end{pmatrix}
= \begin{pmatrix} 1 & 0 & 1 \\ 0 & 1 & -2 \end{pmatrix}.
\]

\[
\fbox{\parbox{0.93\textwidth}{\centering
Beide Wege: $[T]_B^C = \begin{pmatrix} 1 & 0 & 1 \\ 0 & 1 & -2 \end{pmatrix}$.
}}
\]

%% --- 41 (b) ---
\subsubsection*{(b) $[v]_B$ und $[T(v)]_C$ f\"ur $v = (x, y, z)$ und Verifikation}

\textbf{Schritt 1: $[v]_B$.} $(x,y,z) = \alpha b_1 + \beta b_2 + \gamma b_3 = (\alpha + \beta + \gamma,\; \beta + \gamma,\; \gamma)$:
\[
\gamma = z; \quad \beta + \gamma = y \Rightarrow \beta = y - z; \quad \alpha + \beta + \gamma = x \Rightarrow \alpha = x - y.
\]
\[
[v]_B = (x - y,\; y - z,\; z).
\]

\textbf{Schritt 2: $[T(v)]_C$.} Zuerst $T(v) = (x - 2z,\; y - 3z)$. In $C$-Koordinaten ($\beta = q$, $\alpha = p - q$):
\[
\beta = y - 3z; \quad \alpha = (x - 2z) - (y - 3z) = x - y + z.
\]
\[
[T(v)]_C = (x - y + z,\; y - 3z).
\]

\textbf{Schritt 3: Verifikation $[T]_B^C [v]_B = [T(v)]_C$.}
\[
\begin{pmatrix} 1 & 0 & 1 \\ 0 & 1 & -2 \end{pmatrix}\begin{pmatrix} x - y \\ y - z \\ z \end{pmatrix}
= \begin{pmatrix} (x-y) + z \\ (y-z) - 2z \end{pmatrix}
= \begin{pmatrix} x - y + z \\ y - 3z \end{pmatrix}. \;\checkmark
\]

\[
\fbox{\parbox{0.93\textwidth}{\centering
$[v]_B = (x-y,\; y-z,\; z)$, \quad $[T(v)]_C = (x-y+z,\; y-3z)$, \quad Satz best\"atigt. \checkmark
}}
\]

\bigskip
\hrule
\bigskip

%% ============================================================
\subsection*{Aufgabe 42}
%% ============================================================

\textbf{Angabe.} $T \colon \mathbb{R}^{2\times 2} \to \mathbb{R}^{2\times 2}$ mit $T(X) = AX - XA$, wobei $A = \begin{pmatrix} 2 & -1 \\ 1 & 0 \end{pmatrix}$.
Standardbasis $E = \{E_{11}, E_{12}, E_{21}, E_{22}\}$.
\begin{itemize}[nosep]
\item[(a)] $[T]_E$ berechnen.
\item[(b)] Mit $[T]_E$: Basis und Dimension von $\text{Ker}(T)$ und $\text{Im}(T)$.
\end{itemize}

\begin{begriffe}
\begin{itemize}[leftmargin=*, nosep]
    \item $E_{ij}$ ist die Matrix mit einer $1$ an Position $(i,j)$ und sonst Nullen. $\dim \mathbb{R}^{2\times 2} = 4$.
    \item \textbf{Koordinaten einer Matrix} $\begin{pmatrix} a & b \\ c & d \end{pmatrix}$ in $E$: einfach $(a, b, c, d)$ -- die vier Eintr\"age der Reihe nach.
    \item Spalte $j$ von $[T]_E$ $=$ Koordinaten von $T(\text{$j$-ter Basismatrix})$.
\end{itemize}
\end{begriffe}

%% --- 42 (a) ---
\subsubsection*{(a) $[T]_E$ berechnen}

\textbf{Schritt 1: $T$ auf jede Basismatrix anwenden.} Wir brauchen $T(E_{ij}) = A E_{ij} - E_{ij} A$.

\textbf{(a) $T(E_{11})$:}
\[
A E_{11} = \begin{pmatrix} 2 & 0 \\ 1 & 0 \end{pmatrix},\quad
E_{11} A = \begin{pmatrix} 2 & -1 \\ 0 & 0 \end{pmatrix}
\;\Rightarrow\;
T(E_{11}) = \begin{pmatrix} 0 & 1 \\ 1 & 0 \end{pmatrix}
\;\widehat{=}\; (0, 1, 1, 0).
\]

\textbf{(b) $T(E_{12})$:}
\[
A E_{12} = \begin{pmatrix} 0 & 2 \\ 0 & 1 \end{pmatrix},\quad
E_{12} A = \begin{pmatrix} 1 & 0 \\ 0 & 0 \end{pmatrix}
\;\Rightarrow\;
T(E_{12}) = \begin{pmatrix} -1 & 2 \\ 0 & 1 \end{pmatrix}
\;\widehat{=}\; (-1, 2, 0, 1).
\]

\textbf{(c) $T(E_{21})$:}
\[
A E_{21} = \begin{pmatrix} -1 & 0 \\ 0 & 0 \end{pmatrix},\quad
E_{21} A = \begin{pmatrix} 0 & 0 \\ 2 & -1 \end{pmatrix}
\;\Rightarrow\;
T(E_{21}) = \begin{pmatrix} -1 & 0 \\ -2 & 1 \end{pmatrix}
\;\widehat{=}\; (-1, 0, -2, 1).
\]

\textbf{(d) $T(E_{22})$:}
\[
A E_{22} = \begin{pmatrix} 0 & -1 \\ 0 & 0 \end{pmatrix},\quad
E_{22} A = \begin{pmatrix} 0 & 0 \\ 1 & 0 \end{pmatrix}
\;\Rightarrow\;
T(E_{22}) = \begin{pmatrix} 0 & -1 \\ -1 & 0 \end{pmatrix}
\;\widehat{=}\; (0, -1, -1, 0).
\]

\textbf{Schritt 2: Koordinatenvektoren als Spalten.}
\[
[T]_E = \begin{pmatrix}
0 & -1 & -1 & 0 \\
1 & 2 & 0 & -1 \\
1 & 0 & -2 & -1 \\
0 & 1 & 1 & 0
\end{pmatrix}.
\]

%% --- 42 (b) ---
\subsubsection*{(b) Kern und Bild aus $[T]_E$}

\textbf{Schritt 1: $[T]_E$ in Stufenform bringen.}
\[
\begin{pmatrix}
0 & -1 & -1 & 0 \\
1 & 2 & 0 & -1 \\
1 & 0 & -2 & -1 \\
0 & 1 & 1 & 0
\end{pmatrix}
\xrightarrow{Z_1 \leftrightarrow Z_2}
\begin{pmatrix}
1 & 2 & 0 & -1 \\
0 & -1 & -1 & 0 \\
1 & 0 & -2 & -1 \\
0 & 1 & 1 & 0
\end{pmatrix}
\xrightarrow{Z_3 \to Z_3 - Z_1}
\begin{pmatrix}
1 & 2 & 0 & -1 \\
0 & -1 & -1 & 0 \\
0 & -2 & -2 & 0 \\
0 & 1 & 1 & 0
\end{pmatrix}
\]
\[
\xrightarrow[Z_4 \to Z_4 + Z_2]{Z_3 \to Z_3 - 2Z_2}
\begin{pmatrix}
\piv{1} & 2 & 0 & -1 \\
0 & \piv{-1} & -1 & 0 \\
0 & 0 & 0 & 0 \\
0 & 0 & 0 & 0
\end{pmatrix}
\]

Pivots in Spalte 1 und 2 $\Rightarrow$ Rang $= 2$.

\textbf{Schritt 2: Dimensionen.}
\[
\dim \text{Im}(T) = \text{Rang} = 2, \qquad \dim \text{Ker}(T) = 4 - 2 = 2.
\]

\textbf{Schritt 3: Kern bestimmen.} Variablen $(w_1, w_2, w_3, w_4)$ stehen f\"ur $X = w_1 E_{11} + w_2 E_{12} + w_3 E_{21} + w_4 E_{22}$. Die Stufenform liefert:
\begin{align*}
\text{Zeile 2:} &\quad -w_2 - w_3 = 0 \;\Rightarrow\; w_2 = -w_3 \\
\text{Zeile 1:} &\quad w_1 + 2w_2 - w_4 = 0 \;\Rightarrow\; w_1 = -2w_2 + w_4 = 2w_3 + w_4
\end{align*}
Freie Variablen: $w_3, w_4$.
\[
(w_1, w_2, w_3, w_4) = w_3(2, -1, 1, 0) + w_4(1, 0, 0, 1).
\]

Zur\"uck in Matrizen \"ubersetzt:
\[
w_3\text{-Anteil} \;\widehat{=}\; \begin{pmatrix} 2 & -1 \\ 1 & 0 \end{pmatrix} = A, \qquad
w_4\text{-Anteil} \;\widehat{=}\; \begin{pmatrix} 1 & 0 \\ 0 & 1 \end{pmatrix} = I.
\]

\textbf{Das ist anschaulich logisch:} $T(A) = A^2 - A^2 = 0$ und $T(I) = AI - IA = A - A = 0$. Beide liegen wirklich im Kern.

\[
\fbox{\parbox{0.93\textwidth}{\centering
Basis von $\text{Ker}(T)$: $\left\{\begin{pmatrix} 2 & -1 \\ 1 & 0 \end{pmatrix},\; \begin{pmatrix} 1 & 0 \\ 0 & 1 \end{pmatrix}\right\}$. \quad $\dim \text{Ker}(T) = 2$.
}}
\]

\textbf{Schritt 4: Bild bestimmen.} Pivots stehen in Spalte 1 und 2 $\Rightarrow$ die zugeh\"origen Original-Bilder $T(E_{11})$ und $T(E_{12})$ bilden eine Basis des Bildes.

\[
\fbox{\parbox{0.93\textwidth}{\centering
Basis von $\text{Im}(T)$: $\left\{\begin{pmatrix} 0 & 1 \\ 1 & 0 \end{pmatrix},\; \begin{pmatrix} -1 & 2 \\ 0 & 1 \end{pmatrix}\right\}$. \quad $\dim \text{Im}(T) = 2$.
}}
\]

\textbf{Kontrolle:} $\dim \text{Ker} + \dim \text{Im} = 2 + 2 = 4 = \dim \mathbb{R}^{2\times 2}$. \checkmark

\end{document}
